\chapter{Build and deploy}

The repository's root \snapfile{Makefile} is the canonical entry point.
Every target is idempotent and safe to re-run; the build script for this
document (\reposrc{docs/publishing/Makefile}) follows the same convention.

\section{Build chain}

\begin{figure}[!h]
\centering
\snapdiagram{diagrams/build-chain.tex}
\caption{Build chain. Each rectangle is a Make target; arrows are
dependencies. The \texttt{docs-*} targets at the right are the additions
introduced by this document's build pipeline.}
\end{figure}

\subsection{Native (Linux) build}

\begin{lstlisting}[language=bash]
make build-rust       # cargo build --release + (optional) wasm32-unknown-emscripten
make build-godot      # exports every preset in every game project under games/<dir>/
make build-backend    # ./mvnw package -DskipTests
make build-frontend   # ng build for host + four remotes
make build            # runs all of the above in dependency order
\end{lstlisting}

The \texttt{build-rust} target builds the native \snapfile{.so} for the
editor and tries the WASM target; if \texttt{emcc} isn't on \texttt{PATH}
the WASM step prints a warning and continues — the rest of the build
chain doesn't depend on it for normal browser play through Godot's WASM
target.

\subsection{Docker compose (recommended)}

\begin{lstlisting}[language=bash]
docker compose up -d                  # postgres + cassandra + casdoor + bff + frontend + remotes
make e2e                              # full Playwright run against the live stack
make e2e-native                       # same, against the GraalVM native image overlay
\end{lstlisting}

\begin{snapcaution}[title={Docker 29.x needs buildx}]
Docker 29.x removed the legacy in-daemon image builder. \texttt{docker
compose build} will hang silently if \texttt{docker-buildx} isn't
installed. The fix is a user-local install of the v0.18 binary into
\texttt{\textasciitilde/.docker/cli-plugins/docker-buildx} (no sudo
needed). The README's \emph{Running with Docker Compose} section
documents the exact command.
\end{snapcaution}

\section{GraalVM native image}

The BFF compiles to a native binary via GraalVM \texttt{native-image}.
The driver is a separate build profile in \snapfile{backend/pom.xml},
plus a Dockerfile (\snapfile{backend/Dockerfile.native}) and a compose
overlay (\snapfile{compose.native.yml}).

\begin{lstlisting}[language=bash]
make build-backend-native     # mvn -Pnative native:compile (~6 minutes on a workstation)
make e2e-native               # docker compose -f compose.yml -f compose.native.yml up -d
\end{lstlisting}

The native build needs reflection hints registered for every R2DBC
mapping class, every controller DTO, and every Liquibase changelog
entity. \snapfile{R2dbcConvertersConfig} centralises these so adding a
new entity touches one file.

\section{The publishing pipeline}

This document, the Javadoc, and the three slide decks are produced by
\snapfile{docs/publishing/}. The published artefacts are git-ignored.

\begin{lstlisting}[language=bash]
make docs-pdf          # docs/publishing/architecture/build/architecture.pdf
make docs-javadocs     # docs/publishing/javadoc-output/index.html
make docs-slides       # docs/publishing/slides/build/*.pptx
make docs-all          # all of the above
make docs-clean        # remove every produced artefact
\end{lstlisting}

\subsection{LaTeX prerequisites}

\begin{itemize}
  \item TeX Live with \texttt{pdflatex}, \texttt{xelatex},
        \texttt{tcolorbox}, \texttt{tikz}, \texttt{titlesec},
        \texttt{titletoc}, \texttt{enumitem}, \texttt{fontawesome},
        \texttt{inconsolata}.
  \item On Debian / Ubuntu: \texttt{apt install texlive-latex-extra
        texlive-fonts-extra texlive-pictures}.
  \item On Gentoo: the \texttt{tex-live} portage package with USE
        \texttt{extra} produces the same set.
\end{itemize}

\subsection{Javadoc theming}

\snapfile{backend/pom.xml} configures the Maven Javadoc plugin to read
\snapfile{docs/publishing/javadoc-theme/snapdragon.css} as the override
stylesheet and \snapfile{docs/publishing/javadoc-theme/header.html} as
the page header. The override is purely additive — the standard doclet's
default stylesheet still loads first; the Snapdragon CSS layers on top.

\subsection{PowerPoint generator}

The slide decks are generated by Python scripts under
\reposrc{docs/publishing/slides/}. Each script imports
\texttt{python-pptx}, builds the deck programmatically, and writes the
\texttt{.pptx} file to \snapfile{docs/publishing/slides/build/}. The
scripts are deliberately written so a non-Python audience can edit slide
content by editing the data structure at the top of each file —
template placeholders are not used.

\section{CI hooks}

A pre-push hook is shipped at
\reposrc{docs/publishing/scripts/install-hooks.sh}. It installs
\snapfile{.git/hooks/pre-push} that re-runs \texttt{make docs-pdf} on
any push that touches \snapfile{docs/publishing/} or
\snapfile{docs/PLAN-module-federation-and-native.md}. Failures abort the
push so a broken build can't land on \texttt{main}. The hook is opt-in;
the script is idempotent and prints what it changed.
